% Template for post or presentation abstracts submitted to ILAS 2022
% 1. Fill in the presenter's information (name, affiliation, email
% address) and talk information (title of presentation, abstract).
% 2. Compile to make sure there are no errors.
% 3. Save the .tex file for your abstract with a distinctive name,
% ideally "FamilynameGivenname.tex".
% 4. Upload your .tex file to https://clr.ie/129557
%       
% * Please keep your LaTeX commands simple, and avoid the use of
%    user-defined macros.
% * Include details of co-authors and funding acknowledgements after the abstract.
% * Upload only the LaTeX file (with .tex extension).
% * Questions: Contact Niall Madden (Niall.Madden@NUIGalway.ie)

\documentclass[a4paper]{amsart} 
\usepackage{amssymb} 
\usepackage{hyperref}
\pagestyle{empty}
\date{} 

%% IGNORE THE NEXT 4 LINES
\makeatletter % 
\def\labelenumi{\theenumi.}
\def\theenumi{\@arabic\c@enumi}
\makeatother
%% STOP IGNORING NOW

\begin{document}

\title{Geometric continuity, Riordan matrices and applications}
\author{Luis Felipe Prieto-Mart\'inez} %Presenting author only
\address{Universidad Polit\'ecnica de Madrid} % Affiliation of presenting author
\email{luisfelipe.prieto@upm.es} % Email address of presenting author only

\maketitle

% The abstract  ***no more than one page***

%This is a sample abstract for talks and posters
%presented at ILAS 2022, to be held
%in Galway (Ireland) in June 2022. 

%Please, follow the instructions at \url{http://ilas2022.ie/abstracts/}
%for submitting your abstract. Make sure you upload the \texttt{.tex}
%file, and not any others. 

%Try to keep your file as simple as possible, so that it is easily
%rendered in various formats. Give the file a name that readily
%identifies it with you, for example \texttt{OlgaTaussky.tex}. 

%If you wish to include citations, you can do so like this:
%\cite{my_key_for_OT49}. Use a bibitem key that is likely be be unique to
%you (e.g., don't use the default Math Reviews key).
%Any uncited bibliography entries will also
%appear. If you have none, then delete that section.

%% Coauthor details here (optional - delete if not applicable)
% and funding acknowledgment (optional - delete if not applicable)



Geometric continuity $G^k$ is an elementary concept in the geometry of parametrized curves. It is a notion of smoothness that does not depend on a concrete parametrization, but on the curve itself. 

A particular and important question related to this definition is the study of the smoothness of curves described by a piecewise $C^k$ parametrization, that is, the smoothness on the union of  smooth pieces. There are well known compatibility conditions (called the beta-constraints) for the parametrizations of each piece to guarantee such geometric continuity.

In this talk, we will explain how this compatibility conditions can be stated in terms of partial Riordan matrices (for $k<\infty$) and  Riordan matrices (for $k=\infty$). Moreover, we will show how this new statement can help us to prove some uniqueness results concerning analytic curves. We will illustrate this method with two particular results related to well known problems in plane Geometry, namely, (1) if there exists an analytic curve with two interior equichordal points then it must be unique (related to the equichordal problem) and (2) the unique analytic curve with an exterior power point is the circle (related to  Rosenbaum's Power Point Problem).












\medskip







\emph{The author was partially supported by Spanish Goverment grant PGC2018-098321-B-I00.}







\begin{thebibliography}{1}

\bibitem{P}
L. F. Prieto-Mart\'inez.
\newblock Geometric continuity in terms of Riordan matrices and the F-chordal Problem.
\newblock Revista de la Real Academia de Ciencias Exactas, Físicas y Naturales, Seria A, Matemáticas, 115(2):1--15 (2021). 

\bibitem{P2}
L. F. Prieto-Mart\'inez.
\newblock Circular arcs are the only analytic Jordan curves with an exterior power point.
\newblock arXiv preprint arXiv:2201.07892 (2022).


\bibitem{R} M. R. Rychlik.
\newblock  A complete solution to the equichordal point problem of Fujiwara, Blaschke, Rothe and Weitzenb\"ock.
\newblock Inventiones Mathematicae, 129 (1): 141--212 (1997).

\bibitem{Z} L. Zuccheri.
\newblock Characterization of the circle by equipower properties.
\newblock Archiv der Mathematik 58: 199–208  (1992).





%\bibitem{my_key_for_OT49}
%Olga Taussky.
%\newblock A recurring theorem on determinants.
%\newblock {\em  Amer. Math. Monthly}  56:672–676,  (1949).

%\bibitem{another_unique_key}
%Olga Taussky and John Todd.
%\newblock Cholesky, Toeplitz and the triangular factorization of
%symmetric matrices.
%\newblock {\em  Numer. Algorithms}, 41(2):197--202, 2006.
\end{thebibliography}


\end{document}


% Please leave the next bit alone  
%%% Local Variables: 
%%% mode: latex
%%% TeX-master: "../ILAS-2022"
%%% End:
